% T12 C/D, chapitre 7 : Similitudes planes. Leçon et deux derniers exercices : le livre « Terminales scientifiques » (Similitude/Lesona, Similitude/Exo).
% Les autres exercices et tous les corrigés : originaux, écrits en novembre 2026 (les exercices tirés d’autres ouvrages ne sont pas repris, décision du 2 novembre 2026).

\TCDLessonStyleOn
\providecommand{\Plan}{\mathscr{P}}
\providecommand{\Id}{\mathrm{Id}}
\providecommand{\Droite}{(\mathscr{D})}
\providecommand{\DroiteP}{(\mathscr{D}')}
\providecommand{\Arg}[1]{\mathrm{Arg}\left( #1 \right)}
\chapter{Similitudes planes}
\index{Similitude plane}
      \section{G\'en\'eralit\'es}
            Une similitude plane $S$ est une transformation g\'eom\'etrique du plan pour laquelle il existe un r\'eel $k>0$ tel que,
            pour tous points $M$ et $N$ du plan, en posant $M' = S(M)$ et $N' = S(N)$ :
                  \[
                        M'N' = k\, MN.
                  \]
            Le r\'eel $k$ est appel\'e le \emph{rapport} de la similitude $S$.
\index{Similitude plane!rapport}
            Toute isom\'etrie est une similitude de rapport $1$.
            Une homoth\'etie de rapport $k$ ($k \neq 0$) est une similitude de rapport $|k|$.
            %
            \begin{thm}
                  Toute similitude $S$ de rapport $k (>0)$ peut s'\'ecrire sous la forme $S = h \circ f$
                  o\`u $h$ est une homoth\'etie de rapport $k$ et $f$ une isom\'etrie.
                  %
                  \begin{proof}
                        Soit $S$ une similitude de rapport $k$.
                        Soit $h$ l'homoth\'etie de centre $I$ et de rapport $k$.
                        On a $h^{-1} = h\left(I ; \frac{1}{k}\right)$.
                        On pose $f = h^{-1} \circ S$.
                        
                        Montrons que $f$ est une isom\'etrie : \newline
                              Soient $M$ et $N$ deux points, posons $M' = f(M)$ et $N' = f(N)$.
                              On a alors 
                              \begin{align*}
                                    M' & = h^{-1} \left[ S(M) \right]       & N' & = h^{-1} \left[ S(N) \right]      \\
                                       & = h^{-1} (M_1)                     &    & = h^{-1} (N_1)                    \\
                              \end{align*}
                              o\`u $M_1 = S(M)$ et $N_1 = S(N)$, et donc $M_1 N_1 = k MN$.
                              Ainsi, $M'N' = \frac{1}{k} M_1 N_1 = \frac{1}{\not k} \not k MN$.
                              On d\'eduit que $M'N' = MN$, et par cons\'equent, que $f$ est une isom\'etrie.
                              
                        Montrons que $S = h\circ f$ : \newline
                              \[
                                    f = h^{-1} \circ S 
                                          \Rightarrow 
                                                h \circ f = \underbrace{h \circ h^{-1}}_{\Id} \circ S
                                                      \Rightarrow
                                                            h \circ f = S.
                              \]
                  \end{proof}
                  \begin{itemize}
                        \item Si $f$ est un d\'eplacement alors $S$ est une similitude plane directe.
\index{Similitude plane!directe}
                        \item Si $f$ est un antid\'eplacement alors $S$ est une similitude plane indirecte.
\index{Similitude plane!indirecte}
                  \end{itemize}
            \end{thm}
            %
            \begin{prop}
                  \begin{enumerate}
                        \item Toute similitude est une application affine bijective.
                        \item La composition de $2$ similitudes est une similitude.
                              Si
                              \begin{itemize}
                                    \item $S$ est une similitude de rapport $k$,
                                    \item $S'$ une similitude de rapport $k'$,
                              \end{itemize}
                              alors $S \circ S'$ est une similitude de rapport $kk'$.
                        \item Une similitude de rapport $k$ multiplie les distances par $k$ et les aires par $k^2$.
                        \item Une similitude conserve le parall\'elisme, l'orthogonalit\'e, l'alignement et les milieux ;
                              une similitude directe conserve les angles orient\'es, une similitude indirecte les change en leurs oppos\'es.
                  \end{enumerate}
            \end{prop}
            %
      \section{Similitude plane directe}
\index{Similitude plane!directe!forme reduite@forme r\'eduite}
            Une similitude plane directe $S$ est soit une translation
            soit la composition d'une homoth\'etie et d'une rotation de m\^eme centre :
            c'est la forme r\'eduite.
            
            Une similitude plane directe $S$ a pour expression complexe $z' = az + b$ avec $a \in \mathbb{C}^*$ et $b \in \mathbb{C}$.
            \begin{itemize}
                  \item $|a| = k$ est le rapport,
                  \item $\Arg{a} = \alpha$ est l'angle,
                  \item $I\left(\frac{b}{1-a}\right)$ est le centre (si $a\neq 1$).
            \end{itemize}
            
	            Une similitude plane directe est d\'efinie analytiquement par un syst\`eme du type :
	            \[
	                  \left\lbrace
	                  \begin{array}{l}
	                        x' = a x - b y + x_0, \\
	                        y' = b x + a y + y_0.
	                  \end{array}
	                  \right.
	            \]
	            o\`u $a^2 + b^2 = k^2$, avec $k > 0$ le rapport.
      %
      \section{Similitude plane indirecte}
\index{Similitude plane!indirecte}
      %
      \subsection{Forme}
      Une similitude plane indirecte $f$ est la compos\'ee d'une homoth\'etie et d'un antid\'eplacement.
      Le rapport $k$ de $f$ est celui de l'homoth\'etie (en valeur absolue).
      \begin{description}
            \item[Si $k=1$ :] $f$ est un antid\'eplacement.
                  L'ensemble des points invariants de $f$ est soit vide soit une droite.
                  \begin{itemize}
                        \item Si $f$ n'a pas de point invariant, alors $f$ est la compos\'ee d'une sym\'etrie
                              orthogonale et d'une translation de vecteur non nul.
                        \item Si l'ensemble des points invariants est une droite $\Droite$, alors 
                              $f$ la sym\'etrie orthogonale d'axe $\Droite$.
                  \end{itemize}
            \item[Si $k\neq 1$ :] $f$ admet un seul point invariant $\Omega$.
                  Il existe une droite $\Droite$ passant par $\Omega$ telle que 
                  $f = h(\Omega,k) \circ S_{\Droite} = S_{\Droite} \circ h(\Omega,k)$.
                  
                  La droite $\Droite$ est l'ensemble des points $M$ tels que 
                  $\overrightarrow{\Omega M'} = k \cdot \overrightarrow{\Omega M}$,
\index{Similitude plane!indirecte!Premiere droite globalement invariante@Premi\`ere droite globalement invariante}
                  c'est la \emph{premi\`ere droite globalement invariante} par $f$.
                  
                  La droite $\DroiteP = \{ M \mid \overrightarrow{\Omega M'} = -k \cdot \overrightarrow{\Omega M} \}$
\index{Similitude plane!indirecte!Deuxieme droite globalement invariante@Deuxi\`eme droite globalement invariante}
                  est la \emph{deuxi\`eme droite globalement invariante} par $f$.
      \end{description}
%%%%%%%%%%%%%%%%%%%%%%%%%%%%%%%%% Io ambony io ihany no lazaina eto fa toa misy diso (tsy ampy)
%            Une similitude plane indirecte $S$ est la composition d'une homoth\'etie et d'un antid\'eplacement.
%            
%            Un antid\'eplacement est soit une r\'eflexion, 
%            soit une symm\'etrie gliss\'ee ($t_{\vec{u}} \circ S_{\Droite}$ o\`u 
%            $\vec{u}$ est un vecteur directeur de $\Droite$).
%            \begin{itemize}
%                  \item Si $f = S_{\Droite}$, $S = h \circ S_{\Droite}$
%                  \item si $f = t_{\vec{u}} \circ S_{\Droite}$ donc
%                        $S = \underbrace{h \circ t_{\vec{u}} }_{H} \circ S_{\Droite}$ et 
%                        $S = H \circ S_{\Droite}$, o\`u $H$ est une homoth\'etie.
%            \end{itemize}
%            Toute similitude plane indirecte est la composition d'une homoth\'etie et d'une r\'eflexion :
%\index{Similitude plane!indirecte!forme reduite@forme r\'eduite}
%            c'est la forme r\'eduite. \mybox{$S = h \circ S_{\Droite} = S_{\Droite} \circ h$.}
%            
%            Une similitude plane indirecte admet un seul point invariant $I$ : c'est le centre.
%\index{Similitude plane!indirecte!centre}
%\index{Similitude plane!indirecte!point invariant}
%            Il existe une droite $\Droite$ passant par $I$ globalement invariant pas $S$
%            ($M \in \Droite \text{ et } M' = S(M) \Rightarrow \overrightarrow{IM'} = k \overrightarrow{IM}$ ) :
%\index{Similitude plane!indirecte!axe}
%            c'est l'\emph{axe}.
%            
%            \nota{
%                  La droite $(\Delta)$, passant par $I$ et perpandiculaire \`a $\Droite$ 
%                  est aussi globalement invariant par $S$ mais ce n'est pas l'axe : 
%                  c'est la \emph{deuxi\`eme droite globalement invariante par $S$}.
%\index{Similitude plane!indirecte!deuxieme droite globalement invariante@deuxi\`eme droite globalement invariante}
%            }
%      %
%%%%%%%%%%%%%%%%%%%%%%%%%%%%%%%%% 
      \subsection{Expression analytique}
\index{Similitude plane!indirecte!expression analytique}
            Soit $f$ une similitude plane indirecte de rapport $k$.
            La similitude $f$ est alors d\'efinie analytiquement par un syst\`eme du type :
            $\left\lbrace \begin{array}{l} x' = a x + b y + x_0 \\ y' = b x - a y + y_0 \end{array} \right.$
            o\`u $a^2 + b^2 = k^2$.
      %
      \subsection{Expression complexe}
            Soit $S$ la similitude plane indirecte \'egale \`a $h \circ S_{\Droite}$, o\`u $h$ est une homoth\'etie de rapport $k$.
            Soient 
            \[
                  \begin{array}{l}
                        h : z' = k z + b \quad (k \in \mathbb{R}^*) \\
                        S_{\Droite} : z' = a \bar{z} + c 
                  \end{array}
            \]
            les expressions complexes de $h$ et $S_{\Droite}$.
            
            On a alors, $S : z' = k \left( a \bar{z} + c \right) + b = k a \bar{z} + k c + b$.
\index{Similitude plane!indirecte!expression complexe}
            De plus, on v\'erifie que $|k a| = |k| \cdot |a| = |k|$ ($S_{\Droite}$ \'etant une isom\'etrie, $|a| = 1$).
            
            Une similitude plane indirecte a une expression complexe de la forme \thmbox{$z' = a \bar{z} + b$ ($a \neq 0$)}
\index{Similitude plane!indirecte!rapport}
            et son rapport est \thmbox{$k = |a|$}.
            
            \paragraph{Le point invariant $I$}
            Toute similitude plane indirecte de rapport $k \neq 1$ admet un point invariant $I$ qu'on appelle le \emph{centre}.
\index{Similitude plane!indirecte!centre}
                  \[
                        \left.
                              \begin{array}{l}
                                    z_I = a \overline{z_I} + b \\
                                    \overline{z_I} = \bar{a} z_I + \bar{b} 
                              \end{array}
                        \right\rbrace
                              \Rightarrow 
                        z_I = a (\bar{a} z_I + \bar{b} ) + b.
                  \]
            Donc $z_I = a \bar{a} z_I + a \bar{b} + b$.
            Il vient que $(1-a\bar{a}) z_I = a\bar{b} + b$ et donc que
            \thmbox{
                  \[
                        z_I = \frac{a\bar{b}+b}{1 - a \bar{a}} (|a| \neq 1)
                  \]
            }
            est l'affixe du centre $I$.
            %
            \begin{ex}
                  Soit $S$ la transformation g\'eom\'etrique d'expression complexe $z'=2\bar{z}+3i$.
                  Quelle est la nature de $S$\,? D\'eterminer ses \'el\'ements caract\'eristiques.
                  
                  Construire ensuite le transform\'e du cercle $(\mathscr{C})$ de centre $A(1-i)$ et de rayon 1 par $S$.
                  Ecrire l'\'equation de ce cercle.
                  \begin{valiny}
                  %
                  \begin{minipage}{.5\actualtextwidth}
                        \paragraph{Nature et \'el\'ements caract\'eristiques de $S$ :}
                              L'expression est de la forme $z' = a\bar{z} + b$ avec $a = 2$ : $S$ est une similitude plane indirecte,
                              de rapport $|a| = 2 \neq 1$ (ce n'est donc pas un antid\'eplacement).
                              \begin{description}
                                    \item[Rapport :] $k = 2$.
                                    \item[Centre :] $z_I = \frac{-2 \cdot 3i + 3i}{1 - 2 \cdot 2} = \frac{-3i}{-3} = i$.
                                    \item[Axe :] $\Droite = \{ M \mid \overrightarrow{IM'} = 2 \overrightarrow{IM} \}$.
                                          \begin{align*}
                                                \overrightarrow{IM'} & = 2 \overrightarrow{IM} \\
                                                z' - i               & = 2 ( z - i )           \\
                                                2 \bar{z} + 3i - i   & = 2 z - 2 i             \\
                                                2 \bar{z} + 2 i      & = 2 z - 2 i             \\
                                                2 (z - \bar{z})      & = 4 i                   \\
                                                z - \bar{z}          & = 2 i                   \\
                                                x + i y - x + i y    & = 2 i                   \\
                                                2 i y                & = 2 i                   \;.
                                          \end{align*}
                                          Donc $\Droite : y = 1$.
                              \end{description}
                        \paragraph{Transform\'e de $(\mathscr{C})$ :}
                                    $A'(2+5i)$ et $(\mathscr{C}') : (x-2)^2 + (y-5)^2 = 4$.
                  \end{minipage}
                  %
                  \begin{minipage}{.4\actualtextwidth}
                              \resizebox{\linewidth}{!}{%
                              \begin{tikzpicture}[line cap=round,line join=round,>=stealth,x=1cm,y=1cm]
                              \draw [color=black!10, xstep=0.1cm,ystep=0.1cm,line width=.1, opacity=.3] 
                                    (-1,-3) grid (5,8);
                              \draw [color=black!50, xstep=0.5cm,ystep=0.5cm,line width=.2, opacity=.3] 
                                    (-1,-3) grid (5,8);
                              \draw [color=black!100, xstep=1.0cm,ystep=1.0cm,line width=.5, opacity=.3] 
                                    (-1,-3) grid (5,8);
                              \draw[->,color=black] (-1,0) -- (5,0);
                              \foreach \x in {-1,1,2,3,4}
                              \draw[shift={(\x,0)},color=black] (0pt,2pt) -- (0pt,-2pt) node[below] {\footnotesize $\x$};
                              \draw[->,color=black] (0,-3) -- (0,8);
                              \foreach \y in {-3,-2,-1,1,2,3,4,5,6,7}
                              \draw[shift={(0,\y)},color=black] (2pt,0pt) -- (-2pt,0pt) node[left] {\footnotesize $\y$};
                              \draw[color=black] (0pt,-10pt) node[right] {\footnotesize $0$};
                              \clip(-1,-3) rectangle (5,8);
                              \draw[loko1] (1,-1) circle (1);
                              \draw [domain=-1:5,loko2] plot(\x,{(-1-0*\x)/-1});
                              \draw[loko1] (2,5) circle (2);
                              \begin{scriptsize}
                              \fill [color=loko1] (1,-1) circle (1.5pt);
                              \draw[color=loko1] (1.16,-0.74) node {$A$};
                              \draw[color=loko1] (1.54,-0.4) node {$(\mathscr{C})$};
                              \fill [color=black] (0,1) circle (1.5pt);
                              \draw[color=black] (0.1,1.26) node {$I$};
                              \fill [color=loko1] (2,5) circle (1.5pt);
                              \draw[color=loko1] (2.18,5.26) node {$A'$};
                              \draw[color=loko1] (1.58,6.54) node {$(\mathscr{C}')$};
                              \end{scriptsize}
                              \end{tikzpicture}
                              }
                  \end{minipage}
                  %
                  \end{valiny}
            \end{ex}

























\TCDLessonStyleOff
\clearpage
% Similitudes %
\providecommand{\Cercle}{(\mathscr{C})}
\providecommand{\CercleP}{(\mathscr{C}')}
\providecommand{\Droite}{(\mathscr{D})}
\providecommand{\DroiteP}{(\mathscr{D}')}
\providecommand{\Plan}{\mathscr{P}}
\providecommand{\mes}[2]{\mathrm{mes}\left( \widehat{ \overrightarrow{#1} , \overrightarrow{#2} } \right)}
\section{Exercices}
%%%%%%%%%%%%%%%%%%%%%%%%%%%%%%%%%%%%%%%%%%%%%%%%%%%%%%%%%%%%%%%%
\begin{exo}\label{exo:sim-directe-analytique}
      Le plan est rapporté à un repère orthonormé direct $(O;\vec{u},\vec{v})$.
      On considère l'application $f$ qui, au point $M(x;y)$, associe le point $M'(x';y')$ tel que
      \[
            \left\lbrace \begin{array}{l} x' = x + \sqrt{3}\, y - \sqrt{3} \\ y' = -\sqrt{3}\, x + y + \sqrt{3} . \end{array} \right.
      \]
      \begin{enumerate}
            \item Exprimer l'affixe $z'$ de $M'$ en fonction de l'affixe $z$ de $M$.
            \item En déduire que $f$ est une similitude directe dont on donnera le rapport, l'angle et le centre $\Omega$.
            \item Préciser la nature et les éléments caractéristiques de $f^{-1}$.
            \item Déterminer l'image par $f$ :
                  \begin{enumerate}
                        \item du cercle $\Cercle$ de centre $O$ et de rayon $1$ ;
                        \item de l'axe des abscisses (on en donnera une équation).
                  \end{enumerate}
      \end{enumerate}
\end{exo}
%%%%%%%%%%%%%%%%%%%%%%%%%%%%%%%%%%%%%%%%%%%%%%%%%%%%%%%%%%%%%%%%
\begin{exo}\label{exo:sim-reconnaissance}
      Dans le plan complexe rapporté à un repère orthonormé direct, reconnaître chacune des transformations suivantes
      et donner ses éléments caractéristiques (centre, rapport, angle, axe, vecteur selon le cas).
      \begin{enumerate}
            \item $f_1 : z' = -2z + 3i$ ;
            \item $f_2 : z' = (1 - i)z + 2i$ ;
            \item $f_3 : z' = iz + 2$ ;
            \item $f_4 : z' = -i\,\bar{z} + 1 + i$ ;
            \item $f_5 : z' = -\bar{z} + 2 + 2i$ ;
            \item $f_6 : z' = (1 + i\sqrt{3})\,\bar{z}$.
      \end{enumerate}
\end{exo}
%%%%%%%%%%%%%%%%%%%%%%%%%%%%%%%%%%%%%%%%%%%%%%%%%%%%%%%%%%%%%%%%
\begin{exo}\label{exo:sim-deux-points}
      Dans le plan complexe rapporté à un repère orthonormé direct, on donne les points
      $A(2)$, $B(2 + 2i)$, $A'(1)$ et $B'(-1 + 2i)$.
      \begin{enumerate}
            \item Déterminer l'expression complexe de la similitude directe $s$ telle que $s(A) = A'$ et $s(B) = B'$.
            \item Donner le rapport, l'angle et le centre $\Omega$ de $s$.
            \item \begin{enumerate}
                        \item Montrer que $\overrightarrow{AA'}$ s'obtient en tournant $\overrightarrow{\Omega A}$ d'un angle $\frac{\pi}{2}$.
                              En déduire la nature du triangle $\Omega A A'$.
                        \item Expliquer comment construire $\Omega$ à la règle et au compas à partir de $A$ et $A'$ seulement,
                              sachant que $s$ a pour rapport $\sqrt{2}$ et pour angle $\frac{\pi}{4}$.
                  \end{enumerate}
            \item Déterminer l'image par $s$ du cercle de diamètre $[AB]$.
      \end{enumerate}
\end{exo}
%%%%%%%%%%%%%%%%%%%%%%%%%%%%%%%%%%%%%%%%%%%%%%%%%%%%%%%%%%%%%%%%
\begin{exo}\label{exo:sim-composee}
      Dans le plan complexe, on considère le point $A(2)$, l'homothétie $h$ de centre $A$ et de rapport $2$
      et la rotation $r$ de centre $O$ et d'angle $\frac{\pi}{2}$. On pose $s = h \circ r$ et $s' = r \circ h$.
      \begin{enumerate}
            \item Justifier, sans calcul, que $s$ et $s'$ sont des similitudes directes de même rapport et de même angle, que l'on précisera.
            \item Donner les expressions complexes de $s$ et de $s'$ et leurs centres respectifs $\Omega$ et $\Omega'$. A-t-on $s = s'$\,?
            \item \begin{enumerate}
                        \item Montrer, sans calcul, que $s' \circ s^{-1}$ est une translation.
                        \item Calculer $s(O)$ et $s'(O)$ ; en déduire le vecteur de cette translation.
                  \end{enumerate}
      \end{enumerate}
\end{exo}
%%%%%%%%%%%%%%%%%%%%%%%%%%%%%%%%%%%%%%%%%%%%%%%%%%%%%%%%%%%%%%%%
\begin{exo}\label{exo:sim-carre-construction}
      $ABCD$ est un carré de centre $O$ tel que $(\overrightarrow{AB},\overrightarrow{AD}) = \frac{\pi}{2}\ [2\pi]$.
      On note $s$ la similitude directe de centre $A$ qui transforme $B$ en $C$.
      \begin{enumerate}
            \item Donner le rapport et l'angle de $s$.
            \item Construire $D' = s(D)$ et préciser sa position par rapport à $C$ et $D$. Quelle est l'image de $O$\,?
            \item Soit $M$ un point distinct de $A$ et $M' = s(M)$. Montrer que $MM' = MA$ et que
                  $(\overrightarrow{MM'},\overrightarrow{MA}) = \frac{\pi}{2}\ [2\pi]$.
            \item Déterminer l'image par $s$ de la droite $(BD)$.
            \item \emph{Application.} Montrer que, pour tout point $M$ de $(BD)$, il existe un unique point $N$ tel que
                  $MN = MA$ et $(\overrightarrow{MN},\overrightarrow{MA}) = \frac{\pi}{2}$, et que ce point $N$ est sur la droite $(DC)$.
                  Construire $N$ lorsque $M$ est le milieu de $[OB]$.
      \end{enumerate}
\end{exo}
%%%%%%%%%%%%%%%%%%%%%%%%%%%%%%%%%%%%%%%%%%%%%%%%%%%%%%%%%%%%%%%%
\begin{exo}\label{exo:sim-centre-cercles}
      Dans le plan complexe rapporté à un repère orthonormé direct, on donne les points
      $A(0)$, $B(4)$, $A'(2i)$ et $B'(2 + 4i)$. Soit $s$ la similitude directe telle que $s(A) = A'$ et $s(B) = B'$.
      \begin{enumerate}
            \item Déterminer l'expression complexe de $s$, son rapport, son angle et son centre $\Omega$.
            \item Les droites $(AB)$ et $(A'B')$ se coupent en $P$. Déterminer l'affixe de $P$.
            \item \begin{enumerate}
                        \item Justifier que $\left((AB),(A'B')\right) = \frac{\pi}{4}\ [\pi]$ et que $(\overrightarrow{\Omega A},\overrightarrow{\Omega A'}) = \frac{\pi}{4}\ [2\pi]$.
                        \item En déduire que $\Omega$ appartient au cercle passant par $P$, $A$, $A'$, puis au cercle passant par $P$, $B$, $B'$.
                  \end{enumerate}
            \item Déterminer les centres et rayons de ces deux cercles et vérifier par le calcul que leurs points communs sont $P$ et $\Omega$.
                  En déduire une construction géométrique du centre d'une similitude directe connaissant deux points et leurs images.
      \end{enumerate}
\end{exo}
%%%%%%%%%%%%%%%%%%%%%%%%%%%%%%%%%%%%%%%%%%%%%%%%%%%%%%%%%%%%%%%%
\begin{exo}\label{exo:sim-lieu}
      Le plan complexe est rapporté à un repère orthonormé direct $(O;\vec{u},\vec{v})$.
      On note $\Gamma$ le cercle de centre $O$ et de rayon $2$, et $A$ le point d'affixe $4$.
      À tout point $M$ de $\Gamma$, on associe le point $N$ tel que $AN = 2AM$ et $(\overrightarrow{AM},\overrightarrow{AN}) = \frac{\pi}{3}\ [2\pi]$.
      \begin{enumerate}
            \item Soit $K$ le milieu de $[AN]$. Montrer que le triangle $AMK$ est équilatéral, puis que le triangle $AMN$ est rectangle en $M$.
            \item Montrer que $N$ est l'image de $M$ par une similitude directe $s$ dont on donnera le centre, le rapport, l'angle et l'expression complexe.
            \item Déterminer l'ensemble $\Gamma'$ décrit par $N$ lorsque $M$ décrit $\Gamma$.
            \item Exprimer $MN$ en fonction de $AM$. Pour quelles positions de $M$ la longueur $MN$ est-elle minimale\,? maximale\,?
      \end{enumerate}
\end{exo}
%%%%%%%%%%%%%%%%%%%%%%%%%%%%%%%%%%%%%%%%%%%%%%%%%%%%%%%%%%%%%%%%
\begin{exo}\label{exo:sim-spirale}
      Le plan complexe est rapporté à un repère orthonormé direct $(O;\vec{u},\vec{v})$ ; l'unité est le centimètre.
      Soit $s$ la transformation d'expression complexe $z' = \frac{1+i}{2}\, z$.
      On définit les points $M_n$ d'affixe $z_n$ par $M_0(4)$ et $M_{n+1} = s(M_n)$ pour tout entier naturel $n$.
      \begin{enumerate}
            \item Préciser la nature et les éléments caractéristiques de $s$. Placer $M_0$, $M_1$, $M_2$, $M_3$, $M_4$.
            \item Exprimer $z_n$ en fonction de $n$, puis $OM_n$ en fonction de $n$.
            \item Montrer que, pour tout $n$, le triangle $OM_nM_{n+1}$ est rectangle isocèle en $M_{n+1}$.
            \item On note $\ell_n = M_nM_{n+1}$ et $L_n = \ell_0 + \ell_1 + \cdots + \ell_{n-1}$ la longueur de la ligne brisée $M_0M_1\ldots M_n$.
                  Calculer $L_n$ en fonction de $n$ et sa limite quand $n$ tend vers $+\infty$.
            \item Quel est le plus petit entier $n \geq 1$ tel que $M_n$ appartienne à la demi-droite $[O;\vec{u})$\,? Donner alors $z_n$.
      \end{enumerate}
\end{exo}
%%%%%%%%%%%%%%%%%%%%%%%%%%%%%%%%%%%%%%%%%%%%%%%%%%%%%%%%%%%%%%%%
\begin{exo}\label{exo:sim-indirecte}
      Le plan complexe est rapporté à un repère orthonormé direct. Soit $g$ la transformation d'expression complexe
      \[ z' = 2i\,\bar{z} - 1 - i . \]
      \begin{enumerate}
            \item Quelle est la nature de $g$\,? Donner son rapport.
            \item Déterminer l'affixe de son centre $\Omega$.
            \item Déterminer une équation de l'axe $\Droite$ de $g$ (ensemble des points $M$ tels que $\overrightarrow{\Omega M'} = 2\overrightarrow{\Omega M}$)
                  et de la seconde droite $\DroiteP$ globalement invariante (ensemble des points $M$ tels que $\overrightarrow{\Omega M'} = -2\overrightarrow{\Omega M}$).
            \item Vérifier que $g = h \circ s_{\Droite}$, où $h$ est l'homothétie de centre $\Omega$ et de rapport $2$.
            \item Déterminer l'image par $g$ du cercle de centre $E(2)$ et de rayon $1$.
      \end{enumerate}
\end{exo}
%%%%%%%%%%%%%%%%%%%%%%%%%%%%%%%%%%%%%%%%%%%%%%%%%%%%%%%%%%%%%%%%
\begin{exo}\label{exo:sim-indirecte-deux-points}
      Dans le plan complexe rapporté à un repère orthonormé direct, on donne les points
      $A(1)$, $B(1 + i)$, $A'(-1 - i)$ et $B'(-3 - i)$.
      \begin{enumerate}
            \item Montrer qu'il existe une unique similitude indirecte $g$ telle que $g(A) = A'$ et $g(B) = B'$,
                  et donner son expression complexe.
            \item Déterminer le rapport, le centre $\Omega$ et l'axe $\Droite$ de $g$, puis sa forme réduite.
            \item Les droites $(AB)$ et $(A'B')$ sont-elles parallèles\,? perpendiculaires\,? Retrouver ce résultat à l'aide de l'axe $\Droite$.
      \end{enumerate}
\end{exo}
%%%%%%%%%%%%%%%%%%%%%%%%%%%%%%%%%%%%%%%%%%%%%%%%%%%%%%%%%%%%%%%%
\begin{exo}\label{exo:sim-format-a4}
      Une feuille de papier au format A4 a la propriété suivante : en la coupant en deux selon la médiane parallèle au petit côté,
      on obtient deux demi-feuilles de même forme que la feuille entière.
      On modélise la feuille par un rectangle $ABCD$ tel que $AB = \sqrt{2}$, $AD = 1$ et
      $(\overrightarrow{AB},\overrightarrow{AD}) = \frac{\pi}{2}\ [2\pi]$ ;
      on note $I$ et $J$ les milieux respectifs de $[AB]$ et de $[DC]$.
      On rapporte le plan au repère orthonormé direct $(A;\vec{u},\vec{v})$ avec $\vec{u} = \frac{1}{\sqrt{2}}\overrightarrow{AB}$ et $\vec{v} = \overrightarrow{AD}$.
      \begin{description}
            \item[Partie A.] Soit $s$ la similitude directe telle que $s(A) = B$ et $s(B) = C$.
                  \begin{enumerate}
                        \item Déterminer l'expression complexe de $s$, son rapport et son angle.
                        \item Montrer que $s(C) = J$ et $s(D) = I$. Que représente le rectangle $BCJI$\,?
                        \item Déterminer l'affixe du centre $\Omega$ de $s$. Montrer que $\Omega$ est le point d'intersection des droites $(AC)$ et $(BJ)$,
                              et que ces deux droites sont perpendiculaires.
                  \end{enumerate}
            \item[Partie B.] Soit $g$ la similitude indirecte telle que $g(A) = A$ et $g(B) = D$.
                  \begin{enumerate}
                        \item Déterminer l'expression complexe de $g$.
                        \item Calculer $g(C)$ et $g(D)$. Quelle est l'image du rectangle $ABCD$\,?
                        \item Donner le centre, le rapport et l'axe $\Delta$ de $g$, puis sa forme réduite.
                  \end{enumerate}
      \end{description}
\end{exo}
%%%%%%%%%%%%%%%%%%%%%%%%%%%%%%%%%%%%%%%%%%%%%%%%%%%%%%%%%%%%%%%%
\begin{exo}\label{exo:sim-apollonius}
      Le plan complexe est rapporté à un repère orthonormé direct $(O;\vec{u},\vec{v})$.
      On note $A$ et $B$ les points d'affixes $2$ et $1 + i$, et $s$ la similitude directe de centre $O$ qui transforme $A$ en $B$.
      Pour tout point $M$, on note $M' = s(M)$.
      \begin{enumerate}
            \item Déterminer le rapport et l'angle de $s$, puis son expression complexe.
            \item Déterminer l'affixe du point $C$ tel que $s(C) = A$.
            \item Montrer que, pour tout point $M$, $AM' = \frac{\sqrt{2}}{2}\, CM$.
            \item On cherche l'ensemble $\Gamma$ des points $M$ tels que $AM' = OM$.
                  \begin{enumerate}
                        \item Montrer que $M \in \Gamma$ si et seulement si $CM^2 - 2\,OM^2 = 0$.
                        \item Soit $G$ le barycentre de $(C;1)$ et $(O;-2)$. Déterminer l'affixe de $G$ et montrer que
                              $CM^2 - 2\,OM^2 = -MG^2 + 16$.
                        \item En déduire la nature de $\Gamma$. Vérifier que le point $E(2 + 2i)$ appartient à $\Gamma$.
                  \end{enumerate}
      \end{enumerate}
\end{exo}
%%%%%%%%%%%%%%%%%%%%%%%%%%%%%%%%%%%%%%%%%%%%%%%%%%%%%%%%%%%%%%%%
%%%%%%%%%%%%%%%%%%%%%%%%%%%%%%%%%%%%%%%%%%%%%%%%%%%%%%%%%%%%%%%%
\begin{exo}\label{exo:sim-triangle-equilateral}
      \begin{description}
            \item[I.]
                  Dans le plan orient\'e $\Plan$, $a$ \'etant un nombre r\'eel strictement positif,
                  on consid\`ere un triangle \'equilat\'eral $ABC$ de sens direct, de centre $D$, tel que :
                        $\|\overrightarrow{AB}\|=\|\overrightarrow{BC}\|=\|\overrightarrow{CA}\|=a$.
                  \begin{enumerate}
                        \item D\'eterminer le point $G$, barycentre des points $A$, $B$ et $C$
                              respectivement affect\'es des coefficients $2$, $1$ et $1$.
                              Pr\'eciser sur une figure la position du point $G$.
                              (On compl\'etera la figure au fur et \`a mesure.)
                        \item On consid\`ere l'application $f$ de $\Plan$ dans $\Plan$ qui 
                              au point $M$ associe le point $M'$ tel que :
                              $\overrightarrow{MM'} = 2 \overrightarrow{MA}+\overrightarrow{MB}+\overrightarrow{MC}$.
                              Pr\'eciser la nature de $f$ et les \'el\'ements qui la caract\'erisent.
                        \item D\'eterminer l'ensemble $\mathcal{E}$ des points $M$ de $\Plan$ v\'erifiant :
                              $2 \overrightarrow{MA}^2 + \overrightarrow{MB}^2 + \overrightarrow{MC}^2 = 2 a^2$.
                              (On pourra remarquer que $A$ est un \'el\'ement de $\mathcal{E}$.)
                        \item D\'eterminer l'ensemble $\Delta$ des points $M$ de $\Plan$ v\'erifiant :
                              $\overrightarrow{MB}^2 + \overrightarrow{MC}^2 - 2 \overrightarrow{MA}^2 = 2 a^2$.
                              (On pourra remarquer que $A$ est un \'el\'ement de $\Delta$.)
                  \end{enumerate}
            \item[II.] 
                  Dans le plan $\Plan$, soient $I$, $J$ et $K$ les milieux respectifs de
                  $[BC]$, $[CA]$ et $[AB]$.
                  Le point $E$ est le sym\'etrique de $K$ par rapport \`a $I$.
                  \begin{enumerate}
                        \item \begin{enumerate}
                                    \item Montrer qu'il existe une unique rotation $r_1$ qui transforme 
                                          $B$ en $I$ et $J$ en $A$.
                                    \item D\'eterminer les \'el\'ements caract\'eristiques de $r_1$.
                              \end{enumerate}
                        \item Soit $t$ la translation de vecteur $\overrightarrow{AJ}$.
                              On pose : $r_2 = t \circ r_1$ et $f = s_{(JC)} \circ s_{(JE)} \circ s_{(KE)}$.
                              \begin{enumerate}
                                    \item D\'eterminer $r_2 (J)$ et caract\'eriser $r_2$.
                                    \item D\'eterminer deux droites $\Delta_1$ et $\Delta_2$ telles que
                                          $r_1 = s_{(KC)} \circ s_{\Delta_1}$ et $r_2 = s_{(JC)} \circ s_{\Delta_2}$.
                                    \item En d\'eduire que $f = t_{\overrightarrow{AJ}} \circ s_{(KC)}$.
                                    \item D\'eterminer l'image du triangle $BIK$ par $f$.
                                          Justifier que $f$ est une sym\'etrie glissante et d\'eterminer sa forme r\'eduite.
                              \end{enumerate}
                  \end{enumerate}
            \item[III.]
                  Le plan est muni d'un rep\`ere orthonorm\'e direct $(A;\vec{u},\vec{v})$ avec
                  $\vec{u} = \frac{1}{2} \overrightarrow{AB}$, $\|\vec{u}\| = \|\vec{v}\|$ et $a = 4~\mathrm{cm}$.
                  \begin{enumerate}
                        \item \begin{enumerate}
                                    \item D\'eterminer l'\'ecriture complexe de $t_{\overrightarrow{AJ}}$ et $s_{(KC)}$.
                                    \item En d\'eduire l'expression analytique de $f$.
                              \end{enumerate}
                        \item D\'eterminer l'ecriture complexe et les \'el\'ements g\'eom\'etriques
                              de la similitude plane directe $s$ qui transforme $E$ en $I$ et $C$ en $D$.
                  \end{enumerate}
      \end{description}
\end{exo}
%%%%%%%%%%%%%%%%%%%%%%%%%%%%%%%%%%%%%%%%%%%%%%%%%%%%%%%%%%%%%%%%
\begin{exo}\label{exo:sim-triangle-isocele}
      Dans le plan orient\'e $\mathscr{P}$, on consid\`ere un triangle isoc\`ele $ABC$ tel que $AB = AC = a$, 
      o\`u $a$ est un r\'eel strictement positif, et $ \left( \overrightarrow{AB},\overrightarrow{AC} \right) = \frac{\pi}{2} \quad [2 \pi]$.
      On note $D$ le sym\'etrique de $A$ par rapport \`a $B$, $O$ le milieu de $[CD]$ et 
      $(\Gamma)$ le cercle de diam\`etre $[CD]$. 
      Placer sur une figure les points $A$, $B$, $C$, $D$, $O$ et le cercle $(\Gamma)$.       
      %
      \begin{description}
            \item[Partie A.]
                  \begin{enumerate}
                        \item D\'eterminer le barycentre $G$ du syst\`eme de points pond\'er\'es 
                              $\left\lbrace (A,-1) , (B , 2) , (C , 1) \right\rbrace$.
	                        \item D\'eterminer et construire l'ensemble 
	                              \[
	                              \mathscr{C} = \left\lbrace M \in \mathscr{P} \;\middle|\; \left\| - 
	                              \overrightarrow{MA} + 2 \; \overrightarrow{MB} + \overrightarrow{MC} \right\| = a \right\rbrace.
	                              \]
	                        \item D\'eterminer et construire l'ensemble 
	                              \[
	                              \mathscr{C}' = \left\lbrace M \in \mathscr{P} \mid 
	                              -{MA}^2 + 2 \; {MB}^2 + {MC}^2 = a^2 \right\rbrace.
	                              \]
                        \item Montrer que pour tout point $M$ du plan 
                              $\mathscr{P}$, $-\overrightarrow{MA} + 2 \; \overrightarrow{MB} - \overrightarrow{MC}$ 
	                              est un vecteur constant $\overrightarrow{W}$ que l'on pr\'ecisera.
                  \end{enumerate}
            \item[Partie B]
                  On d\'esigne par $s$ la similitude directe qui transforme $D$ en $B$ et $B$ en $C$. 
                  On se propose de d\'eterminer, par deux mani\`eres ind\'ependantes, 
                  les \'el\'ements caract\'eristiques de $s$, notamment son centre $I$.
                  
                  \textbf{I. M\'ethode g\'eom\'etrique}
                  
                  \begin{enumerate}
                        \item D\'eterminer le rapport $k$ et l'angle $\alpha$ de la similitude $s$. 
                        \item Montrer que : 
                              \begin{equation} \label{eq:angle}
                                    (\overrightarrow{ID}, \overrightarrow{IC}) = -\frac{\pi}{2} \quad [2 \pi] \;,
                              \end{equation}
                              %
                              et
                              %
                              \begin{equation} \label{eq:rapport}
                                    IC = 2 \; ID \;.
                              \end{equation}
                        \item A l'aide de (\ref{eq:angle}), prouver que $I$ appartient au cercle $(\Gamma)$ ; 
                              puis, en utilisant (\ref{eq:rapport}), que $ID = a$.
                              Etablir enfin que $BI = BC$.
                        \item Prouver que la droite $(OB)$ est la m\'ediatrice de $[IC]$.
                              Pr\'eciser la nature du quadrilat\`ere $CADI$. Placer le point $I$.
                  \end{enumerate}
                  
                  \textbf{II. Utilisation des nombres complexes}
                  
                  On pose $\vec{u} = \frac{1}{a} \overrightarrow{AB}$ et on rapporte le 
                  plan complexe au rep\`ere orthonormal $(A, \vec{u}, \vec{v})$.
                  On note $z_0$ l'affixe de $I$.
                  \begin{enumerate}
                        \item D\'eterminer les affixes des points $B$, $C$ et $D$.
                        \item D\'eterminer l'\'ecriture complexe de la similitude $s$.
                              D\'eterminer $z_0$ et pr\'eciser la position de $I$.
                  \end{enumerate}
                  
                  \textbf{III.} 
                        Soit $f$ la similitude plane indirecte de centre $A$ qui transforme $C$ en $D$.       
                        D\'eterminer l'\'ecriture complexe de $f$ et donner ses \'el\'ements caract\'eristiques.
      \end{description}
\end{exo}

\begin{corrige}
\section{Solutions des exercices}
%%%%%%%%%%%%%%%%%%%%%%%%%%%%%%%%%%%%%%%%%%%%%%%%%%%%%%%%%%%%%%%%
\begin{sol}[Exercice~\ref{exo:sim-directe-analytique}]
\begin{enumerate}
      \item $z' = x' + iy' = x + \sqrt3 y - \sqrt3 + i(-\sqrt3 x + y + \sqrt3)$.
            Or $(1 - i\sqrt3)(x + iy) = x + \sqrt3 y + i(y - \sqrt3 x)$. Donc
            \fbox{$z' = (1 - i\sqrt3)\,z - \sqrt3 + i\sqrt3$.}
      \item $f$ est de la forme $z' = az + b$ avec $a = 1 - i\sqrt3 = 2e^{-i\pi/3}$ : c'est une similitude directe de rapport $|a| = 2$ et d'angle $-\frac{\pi}{3}$.
            Comme $a \neq 1$, son centre a pour affixe
            $\frac{b}{1-a} = \frac{-\sqrt3 + i\sqrt3}{i\sqrt3} = \frac{-1 + i}{i} = 1 + i$. \fbox{$\Omega(1+i)$, rapport $2$, angle $-\frac{\pi}{3}$.}
      \item $f^{-1}$ est la similitude directe de même centre $\Omega$, de rapport $\frac12$ et d'angle $\frac{\pi}{3}$.
      \item \begin{enumerate}
                  \item Une similitude de rapport $2$ transforme un cercle de rayon $R$ en un cercle de rayon $2R$, de centre l'image du centre.
                        $f(O)$ a pour affixe $b = -\sqrt3 + i\sqrt3$.
                        \fbox{$f(\Cercle)$ : cercle de centre $(-\sqrt3;\sqrt3)$ et de rayon $2$.}
                  \item Pour $M(x;0)$ : $x' = x - \sqrt3$ et $y' = -\sqrt3 x + \sqrt3$. En remplaçant $x = x' + \sqrt3$ :
                        $y' = -\sqrt3 x' - 3 + \sqrt3$. \fbox{L'image est la droite d'équation $y = -\sqrt3\,x + \sqrt3 - 3$.}
                        (Elle fait avec l'axe des abscisses un angle de $-\frac{\pi}{3}$, l'angle de $f$.)
            \end{enumerate}
\end{enumerate}
\end{sol}
%%%%%%%%%%%%%%%%%%%%%%%%%%%%%%%%%%%%%%%%%%%%%%%%%%%%%%%%%%%%%%%%
\begin{sol}[Exercice~\ref{exo:sim-reconnaissance}]
On rappelle : $z' = az + b$ (similitude directe ; centre $\frac{b}{1-a}$ si $a \neq 1$) et $z' = a\bar{z} + b$ (similitude indirecte, antidéplacement si $|a| = 1$).
\begin{enumerate}
      \item $a = -2$ réel : homothétie de rapport $-2$ et de centre d'affixe $\frac{3i}{1+2} = i$.
      \item $a = 1 - i = \sqrt2\, e^{-i\pi/4}$ : similitude directe de rapport $\sqrt2$, d'angle $-\frac{\pi}{4}$ et de centre d'affixe $\frac{2i}{1 - (1-i)} = \frac{2i}{i} = 2$.
      \item $a = i$, $|a| = 1$ : rotation d'angle $\frac{\pi}{2}$ et de centre d'affixe $\frac{2}{1-i} = 1 + i$.
      \item $a = -i$, $b = 1 + i$, $|a| = 1$ ; $a\bar{b} + b = -i(1 - i) + 1 + i = -i - 1 + 1 + i = 0$ : réflexion.
            Points invariants : $x + iy = -i(x - iy) + 1 + i = -y + 1 + i(1 - x)$, soit $x + y = 1$. Axe : la droite $x + y = 1$.
      \item $a = -1$, $b = 2 + 2i$ ; $a\bar{b} + b = -(2 - 2i) + 2 + 2i = 4i \neq 0$ : symétrie glissée.
            $x' = -x + 2$, $y' = y + 2$ : c'est la réflexion d'axe $x = 1$ suivie de la translation de vecteur $(0;2)$, qui dirige cet axe.
      \item $a = 1 + i\sqrt3 = 2e^{i\pi/3}$, $|a| = 2 \neq 1$ : similitude indirecte de rapport $2$ ; son centre est $O$ ($b = 0$).
            Axe : $M$ tel que $z' = 2z$, soit $e^{i\pi/3}\bar{z} = z$, c'est-à-dire $\arg z \equiv \frac{\pi}{6}\ [\pi]$ (ou $z = 0$) :
            la droite passant par $O$ d'équation $y = \frac{x}{\sqrt3}$. Ainsi $f_6 = h(O,2) \circ s_{\Delta}$ avec $\Delta : y = \frac{x}{\sqrt3}$.
\end{enumerate}
\end{sol}
%%%%%%%%%%%%%%%%%%%%%%%%%%%%%%%%%%%%%%%%%%%%%%%%%%%%%%%%%%%%%%%%
\begin{sol}[Exercice~\ref{exo:sim-deux-points}]
\begin{enumerate}
      \item $s : z' = az + b$ avec $a = \frac{z_{B'} - z_{A'}}{z_B - z_A} = \frac{-2 + 2i}{2i} = \frac{(-2+2i)(-i)}{2} = 1 + i$
            et $b = z_{A'} - a z_A = 1 - 2(1+i) = -1 - 2i$. \fbox{$s : z' = (1+i)z - 1 - 2i$.}
      \item $1 + i = \sqrt2\, e^{i\pi/4}$ : rapport $\sqrt2$, angle $\frac{\pi}{4}$ ;
            centre d'affixe $\frac{-1-2i}{1 - (1+i)} = \frac{-1 - 2i}{-i} = 2 - i$. \fbox{$\Omega(2 - i)$.}
      \item \begin{enumerate}
                  \item De $z_{A'} - \omega = (1+i)(z_A - \omega)$ on tire $z_{A'} - z_A = i(z_A - \omega)$ :
                        $\overrightarrow{AA'}$ est l'image de $\overrightarrow{\Omega A}$ par la rotation vectorielle d'angle $\frac{\pi}{2}$.
                        (Ici $z_{A'} - z_A = -1$ et $i(z_A - \omega) = i \times i = -1$.)
                        Donc $AA' = \Omega A$ et $(\Omega A) \perp (AA')$ : \fbox{le triangle $\Omega AA'$ est rectangle isocèle en $A$.}
                  \item On trace la perpendiculaire à $(AA')$ en $A$ et on y reporte, au compas, la longueur $AA'$ : on obtient deux points.
                        $\Omega$ est celui pour lequel $(\overrightarrow{\Omega A},\overrightarrow{\Omega A'}) = +\frac{\pi}{4}$
                        (le triangle $\Omega A A'$ rectangle isocèle en $A$ est alors direct dans l'ordre $\Omega$, $A$, $A'$).
            \end{enumerate}
      \item Le cercle de diamètre $[AB]$ a pour centre le milieu de $[AB]$, d'affixe $2 + i$, et pour rayon $1$.
            $s(2+i) = (1+i)(2+i) - 1 - 2i = 1 + 3i - 1 - 2i = i$.
            \fbox{Image : le cercle de centre $(0;1)$ et de rayon $\sqrt2$} (c'est aussi le cercle de diamètre $[A'B']$).
\end{enumerate}
\end{sol}
%%%%%%%%%%%%%%%%%%%%%%%%%%%%%%%%%%%%%%%%%%%%%%%%%%%%%%%%%%%%%%%%
\begin{sol}[Exercice~\ref{exo:sim-composee}]
\begin{enumerate}
      \item $h$ est une similitude directe de rapport $2$ et d'angle $0$, $r$ une similitude directe de rapport $1$ et d'angle $\frac{\pi}{2}$.
            Les rapports se multiplient et les angles s'ajoutent : $s$ et $s'$ sont des similitudes directes de rapport $2$ et d'angle $\frac{\pi}{2}$.
      \item $r : z' = iz$ et $h : z' = 2(z - 2) + 2 = 2z - 2$.
            \begin{itemize}
                  \item $s(z) = h(iz) = 2iz - 2$ ; centre $\frac{-2}{1 - 2i} = \frac{-2(1+2i)}{5} = -\frac25 - \frac45 i$.
                  \item $s'(z) = r(2z - 2) = 2iz - 2i$ ; centre $\frac{-2i}{1-2i} = \frac{-2i(1+2i)}{5} = \frac45 - \frac25 i$.
            \end{itemize}
            Les centres sont différents : \fbox{$s \neq s'$} (la composition n'est pas commutative).
      \item \begin{enumerate}
                  \item $s^{-1}$ a pour rapport $\frac12$ et pour angle $-\frac{\pi}{2}$ ; $s' \circ s^{-1}$ a donc pour rapport $1$ et pour angle $0$ :
                        c'est un déplacement d'angle nul, une translation.
                  \item $s(O) = h(r(O)) = h(O)$, d'affixe $-2$ ; $s'(O) = r(h(O)) = r(-2)$, d'affixe $-2i$.
                        Donc $(s' \circ s^{-1})(s(O)) = s'(O)$ : la translation envoie le point d'affixe $-2$ sur le point d'affixe $-2i$.
                        \fbox{Vecteur d'affixe $-2i - (-2) = 2 - 2i$.}
                        (Vérification : $s^{-1}(z) = \frac{z+2}{2i}$ et $s'(s^{-1}(z)) = z + 2 - 2i$.)
            \end{enumerate}
\end{enumerate}
\end{sol}
%%%%%%%%%%%%%%%%%%%%%%%%%%%%%%%%%%%%%%%%%%%%%%%%%%%%%%%%%%%%%%%%
\begin{sol}[Exercice~\ref{exo:sim-carre-construction}]
On utilise le repère orthonormé direct $(A;\overrightarrow{AB},\overrightarrow{AD})$ : $z_A = 0$, $z_B = 1$, $z_C = 1 + i$, $z_D = i$, $z_O = \frac{1+i}{2}$.
\begin{enumerate}
      \item Rapport $\frac{AC}{AB} = \sqrt2$, angle $(\overrightarrow{AB},\overrightarrow{AC}) = \frac{\pi}{4}$. Ainsi $s : z' = (1+i)z$.
      \item $z_{D'} = (1+i)i = -1 + i = 2z_D - z_C$ : \fbox{$D'$ est le symétrique de $C$ par rapport à $D$.}
            $s(O)$ a pour affixe $(1+i)\frac{1+i}{2} = i$ : \fbox{$s(O) = D$.}
      \item $z_{M'} - z_M = (1+i)z_M - z_M = i z_M$ : le vecteur $\overrightarrow{MM'}$ s'obtient en tournant $\overrightarrow{AM}$ de $\frac{\pi}{2}$.
            Donc $MM' = AM$, et $\frac{z_A - z_M}{z_{M'} - z_M} = \frac{-z_M}{iz_M} = i$, d'où $(\overrightarrow{MM'},\overrightarrow{MA}) = \frac{\pi}{2}$.
      \item $s(B) = C$ et $s(D) = D'$ ; une similitude transforme une droite en droite : $s((BD)) = (CD')$, qui est la droite $(DC)$
            puisque $D' \in (DC)$.
      \item La condition $MN = MA$, $(\overrightarrow{MN},\overrightarrow{MA}) = \frac{\pi}{2}$ signifie que $N$ est l'image de $A$ par la rotation de centre $M$ et d'angle $-\frac{\pi}{2}$ :
            $N$ existe et est unique. D'après 3, $s(M)$ vérifie cette condition, donc $N = s(M)$, et $N \in s((BD)) = (DC)$.
            Pour $M$ milieu de $[OB]$ : $z_M = \frac{3+i}{4}$ et $z_N = (1+i)\frac{3+i}{4} = \frac{2 + 4i}{4} = \frac12 + i$ :
            \fbox{$N$ est le milieu de $[DC]$.}
\end{enumerate}
\end{sol}
%%%%%%%%%%%%%%%%%%%%%%%%%%%%%%%%%%%%%%%%%%%%%%%%%%%%%%%%%%%%%%%%
\begin{sol}[Exercice~\ref{exo:sim-centre-cercles}]
\begin{enumerate}
      \item $a = \frac{(2+4i) - 2i}{4 - 0} = \frac{1+i}{2}$ et $b = z_{A'} - a z_A = 2i$ : \fbox{$s : z' = \frac{1+i}{2}z + 2i$.}
            Rapport $\left|\frac{1+i}{2}\right| = \frac{\sqrt2}{2}$, angle $\frac{\pi}{4}$, centre d'affixe
            $\frac{2i}{1 - \frac{1+i}{2}} = \frac{4i}{1-i} = 2i(1+i) = -2 + 2i$. \fbox{$\Omega(-2+2i)$.}
      \item $(AB)$ est l'axe des abscisses ; $(A'B')$ passe par $(0;2)$ avec le vecteur directeur $(2;2)$ : $y = x + 2$. \fbox{$P(-2;0)$.}
      \item \begin{enumerate}
                  \item $s$ transforme la droite $(AB)$ en $(A'B')$ et une similitude directe d'angle $\theta$ fait tourner les droites de $\theta$ :
                        $((AB),(A'B')) = \frac{\pi}{4}\ [\pi]$ (ce qu'on voit aussi sur les vecteurs directeurs $(1;0)$ et $(1;1)$).
                        Comme $s(A) = A'$, $\Omega A' = \frac{\sqrt2}{2}\Omega A$ et $(\overrightarrow{\Omega A},\overrightarrow{\Omega A'}) = \frac{\pi}{4}\ [2\pi]$.
                  \item Les droites $(PA)$ et $(PA')$ sont $(AB)$ et $(A'B')$, donc $((PA),(PA')) = ((\Omega A),(\Omega A')) = \frac{\pi}{4}\ [\pi]$ :
                        les points $\Omega$, $P$, $A$, $A'$ sont cocycliques (théorème de l'angle inscrit, version angles de droites).
                        De même avec $B$ et $B'$ : $((PB),(PB')) = ((\Omega B),(\Omega B')) = \frac{\pi}{4}\ [\pi]$.
            \end{enumerate}
      \item Cercle $(PAA')$ : centre sur la médiatrice $x = -1$ de $[PA]$ et sur la médiatrice $y = 1$ de $[AA']$ : centre $(-1;1)$, rayon $\sqrt2$.
            Cercle $(PBB')$ : centre sur $x = 1$ (médiatrice de $[PB]$) et équidistant de $P$ et $B'$ :
            $9 + y^2 = 1 + (y-4)^2$ donne $y = 1$ ; centre $(1;1)$, rayon $\sqrt{10}$.
            Points communs : $(x+1)^2 + (y-1)^2 = 2$ et $(x-1)^2 + (y-1)^2 = 10$ ; par différence $4x = -8$, $x = -2$, puis $(y-1)^2 = 1$ :
            $y = 0$ ou $y = 2$. \fbox{Les points communs sont $P(-2;0)$ et $\Omega(-2;2)$.}

            Construction : si $(AB)$ et $(A'B')$ se coupent en $P$, le centre de la similitude directe envoyant $A$ sur $A'$ et $B$ sur $B'$
            est le second point d'intersection des cercles circonscrits aux triangles $PAA'$ et $PBB'$.
\end{enumerate}
\end{sol}
%%%%%%%%%%%%%%%%%%%%%%%%%%%%%%%%%%%%%%%%%%%%%%%%%%%%%%%%%%%%%%%%
\begin{sol}[Exercice~\ref{exo:sim-lieu}]
\begin{enumerate}
      \item $AK = \frac12 AN = AM$ et $(\overrightarrow{AM},\overrightarrow{AK}) = (\overrightarrow{AM},\overrightarrow{AN}) = \frac{\pi}{3}$ :
            le triangle $AMK$ est isocèle en $A$ avec un angle au sommet de $\frac{\pi}{3}$, donc équilatéral.
            Alors $KM = KA = KN$ : $M$ est sur le cercle de diamètre $[AN]$, donc \fbox{$AMN$ est rectangle en $M$.}
      \item $AN = 2AM$ et $(\overrightarrow{AM},\overrightarrow{AN}) = \frac{\pi}{3}$ : $N = s(M)$ où $s$ est la similitude directe
            de centre $A$, de rapport $2$ et d'angle $\frac{\pi}{3}$.
            \fbox{$s : z' - 4 = 2e^{i\pi/3}(z - 4)$, soit $z' = (1 + i\sqrt3)z - 4i\sqrt3$.}
      \item $\Gamma' = s(\Gamma)$ est le cercle de centre $s(O)$, d'affixe $-4i\sqrt3$, et de rayon $2 \times 2 = 4$.
      \item Dans le triangle $AMN$ rectangle en $M$ : $MN^2 = AN^2 - AM^2 = 3AM^2$, donc \fbox{$MN = \sqrt3\, AM$.}
            Pour $M \in \Gamma$, $AM$ varie entre $OA - 2 = 2$ (pour $M(2)$) et $OA + 2 = 6$ (pour $M(-2)$).
            $MN$ est minimale, égale à $2\sqrt3$, pour $M(2)$, et maximale, égale à $6\sqrt3$, pour $M(-2)$.
\end{enumerate}
\end{sol}
%%%%%%%%%%%%%%%%%%%%%%%%%%%%%%%%%%%%%%%%%%%%%%%%%%%%%%%%%%%%%%%%
\begin{sol}[Exercice~\ref{exo:sim-spirale}]
\begin{enumerate}
      \item $\frac{1+i}{2} = \frac{\sqrt2}{2}e^{i\pi/4}$ : $s$ est la similitude directe de centre $O$, de rapport $\frac{\sqrt2}{2}$ et d'angle $\frac{\pi}{4}$.
            $z_1 = 2 + 2i$, $z_2 = 2i$, $z_3 = -1 + i$, $z_4 = -1$.
      \item $(z_n)$ est géométrique de raison $\frac{1+i}{2}$ : \fbox{$z_n = 4\left(\frac{1+i}{2}\right)^n = 4\left(\frac{\sqrt2}{2}\right)^n e^{in\pi/4}$}
            et $OM_n = 4\left(\frac{\sqrt2}{2}\right)^n$.
      \item $\frac{z_n - z_{n+1}}{0 - z_{n+1}} = \frac{z_n\left(1 - \frac{1+i}{2}\right)}{-\frac{1+i}{2}z_n} = -\frac{1 - i}{1 + i} = -(-i) = i$.
            Donc $M_{n+1}M_n = M_{n+1}O$ et $(\overrightarrow{M_{n+1}O},\overrightarrow{M_{n+1}M_n}) = \frac{\pi}{2}$ :
            le triangle $OM_nM_{n+1}$ est rectangle isocèle en $M_{n+1}$.
      \item D'après 3, $\ell_n = OM_{n+1} = 4\left(\frac{\sqrt2}{2}\right)^{n+1}$ : suite géométrique de premier terme $\ell_0 = 2\sqrt2$ et de raison $q = \frac{\sqrt2}{2}$.
            \[ L_n = 2\sqrt2 \cdot \frac{1 - q^n}{1 - q} = \frac{4\sqrt2}{2 - \sqrt2}\left(1 - \left(\tfrac{\sqrt2}{2}\right)^n\right) = (4 + 4\sqrt2)\left(1 - \left(\tfrac{\sqrt2}{2}\right)^n\right). \]
            Comme $0 < q < 1$, $q^n \to 0$ : \fbox{$\lim L_n = 4 + 4\sqrt2 \approx 9{,}66~\mathrm{cm}$.}
      \item $M_n \in [O;\vec{u})$ si et seulement si $\frac{n\pi}{4} \equiv 0\ [2\pi]$, c'est-à-dire $n$ multiple de $8$.
            \fbox{$n = 8$ et $z_8 = 4 \times \left(\frac{\sqrt2}{2}\right)^8 = \frac{4}{16} = \frac14$.}
\end{enumerate}
\end{sol}
%%%%%%%%%%%%%%%%%%%%%%%%%%%%%%%%%%%%%%%%%%%%%%%%%%%%%%%%%%%%%%%%
\begin{sol}[Exercice~\ref{exo:sim-indirecte}]
\begin{enumerate}
      \item $g$ est de la forme $z' = a\bar{z} + b$ avec $a = 2i$ : similitude indirecte de rapport $|a| = 2$.
      \item $z_\Omega = \frac{a\bar{b} + b}{1 - a\bar{a}} = \frac{2i(-1 + i) - 1 - i}{1 - 4} = \frac{-3 - 3i}{-3} = 1 + i$.
            (Vérification : $2i(1 - i) - 1 - i = 2i + 2 - 1 - i = 1 + i$.)
      \item $\overrightarrow{\Omega M'} = 2\overrightarrow{\Omega M}$ s'écrit $2i\bar{z} - 1 - i - (1+i) = 2(z - 1 - i)$, soit $i\bar{z} = z$.
            Avec $z = x + iy$ : $y + ix = x + iy$, d'où \fbox{$\Droite : y = x$.}
            $\overrightarrow{\Omega M'} = -2\overrightarrow{\Omega M}$ s'écrit $2i\bar{z} - 2 - 2i = -2z + 2 + 2i$, soit $i\bar{z} + z = 2 + 2i$,
            c'est-à-dire $x + y = 2$ (deux fois). \fbox{$\DroiteP : y = -x + 2$}, perpendiculaire à $\Droite$ en $\Omega$.
      \item $s_{\Droite} : z' = i\bar{z}$ (réflexion d'axe $y = x$) et $h : z' = 2(z - 1 - i) + 1 + i = 2z - 1 - i$.
            $h \circ s_{\Droite} : z' = 2i\bar{z} - 1 - i$ : c'est bien $g$.
      \item $g(2) = 4i - 1 - i = -1 + 3i$ ; \fbox{image : cercle de centre $(-1;3)$ et de rayon $2$.}
\end{enumerate}
\end{sol}
%%%%%%%%%%%%%%%%%%%%%%%%%%%%%%%%%%%%%%%%%%%%%%%%%%%%%%%%%%%%%%%%
\begin{sol}[Exercice~\ref{exo:sim-indirecte-deux-points}]
\begin{enumerate}
      \item On cherche $g : z' = a\bar{z} + b$ ($a \neq 0$) avec $a\overline{z_A} + b = z_{A'}$ et $a\overline{z_B} + b = z_{B'}$.
            Par différence, $a\,\overline{(z_B - z_A)} = z_{B'} - z_{A'}$, soit $a(-i) = -2$, d'où $a = -2i$ (unique), puis
            $b = z_{A'} - a\overline{z_A} = -1 - i + 2i = -1 + i$. \fbox{$g : z' = -2i\,\bar{z} - 1 + i$.}
      \item Rapport $|a| = 2$. Centre : $\frac{a\bar{b} + b}{1 - |a|^2} = \frac{-2i(-1 - i) - 1 + i}{-3} = \frac{2i - 2 - 1 + i}{-3} = 1 - i$.
            Axe : $g(z) - z_\Omega = 2(z - z_\Omega)$ s'écrit $-2i\bar{z} - 1 + i - (1 - i) = 2z - 2(1 - i)$, soit $-2i\bar{z} = 2z$, c'est-à-dire $-i\bar{z} = z$ ;
            avec $z = x + iy$ : $-y - ix = x + iy$, d'où \fbox{$\Droite : y = -x$} (qui passe bien par $\Omega(1;-1)$).
            Forme réduite : \fbox{$g = h(\Omega, 2) \circ s_{\Droite}$}. (Vérification : $s_{\Droite} : z' = -i\bar{z}$, $h : z' = 2z - 1 + i$.)
      \item $\overrightarrow{AB}$ a pour affixe $i$ (vertical) et $\overrightarrow{A'B'}$ pour affixe $-2$ (horizontal) : les droites sont perpendiculaires.
            Une homothétie conserve les directions ; la réflexion d'axe $\Droite$, qui fait l'angle $-\frac{\pi}{4}$ avec $(O;\vec{u})$, envoie une droite d'angle $\frac{\pi}{2}$
            sur une droite d'angle $2 \times (-\frac{\pi}{4}) - \frac{\pi}{2} = -\pi$, c'est-à-dire horizontale : on retrouve $(A'B') \perp (AB)$.
\end{enumerate}
\end{sol}
%%%%%%%%%%%%%%%%%%%%%%%%%%%%%%%%%%%%%%%%%%%%%%%%%%%%%%%%%%%%%%%%
\begin{sol}[Exercice~\ref{exo:sim-format-a4}]
Affixes : $z_A = 0$, $z_B = \sqrt2$, $z_C = \sqrt2 + i$, $z_D = i$, $z_I = \frac{\sqrt2}{2}$, $z_J = \frac{\sqrt2}{2} + i$.
\begin{description}
      \item[Partie A.]
      \begin{enumerate}
            \item $a = \frac{z_C - z_B}{z_B - z_A} = \frac{i}{\sqrt2} = \frac{\sqrt2}{2}i$, $b = z_B - a z_A = \sqrt2$.
                  \fbox{$s : z' = \frac{\sqrt2}{2}i\,z + \sqrt2$}, de rapport $\frac{\sqrt2}{2}$ et d'angle $\frac{\pi}{2}$.
            \item $s(C) = \frac{\sqrt2}{2}i(\sqrt2 + i) + \sqrt2 = i - \frac{\sqrt2}{2} + \sqrt2 = \frac{\sqrt2}{2} + i = z_J$ et
                  $s(D) = \frac{\sqrt2}{2}i \cdot i + \sqrt2 = \frac{\sqrt2}{2} = z_I$.
                  $s$ transforme le rectangle $ABCD$ en le rectangle $BCJI$, qui est la demi-feuille : la demi-feuille est bien une réduction
                  de la feuille entière, de rapport $\frac{\sqrt2}{2}$.
            \item $z_\Omega = \frac{\sqrt2}{1 - \frac{\sqrt2}{2}i} = \frac{2}{\sqrt2 - i} = \frac{2(\sqrt2 + i)}{3}$.
                  \fbox{$\Omega\left(\frac{2\sqrt2}{3};\frac23\right)$.}
                  $z_\Omega = \frac23 z_C$, donc $\Omega \in (AC)$ ; $z_\Omega - z_B = \frac{-\sqrt2 + 2i}{3} = \frac23(z_J - z_B)$, donc $\Omega \in (BJ)$.
                  Les vecteurs $\overrightarrow{AC}(\sqrt2;1)$ et $\overrightarrow{BJ}\left(-\frac{\sqrt2}{2};1\right)$ ont pour produit scalaire $-1 + 1 = 0$ :
                  \fbox{$(AC) \perp (BJ)$ en $\Omega$.}
                  (C'est aussi une conséquence de $s((AC)) = (BJ)$ et de l'angle $\frac{\pi}{2}$ de $s$.)
      \end{enumerate}
      \item[Partie B.]
      \begin{enumerate}
            \item $g(A) = A$ donne $b = 0$, donc $g : z' = a\bar{z}$ ; $g(B) = D$ donne $a\sqrt2 = i$, $a = \frac{\sqrt2}{2}i$.
                  \fbox{$g : z' = \frac{\sqrt2}{2}i\,\bar{z}$.}
            \item $g(C) = \frac{\sqrt2}{2}i(\sqrt2 - i) = i + \frac{\sqrt2}{2} = z_J$ et $g(D) = \frac{\sqrt2}{2}i(-i) = \frac{\sqrt2}{2} = z_I$.
                  L'image de $ABCD$ est le rectangle $ADJI$ (l'autre demi-feuille).
            \item Centre $A$ (seul point invariant, car le rapport est différent de $1$), rapport $\frac{\sqrt2}{2}$.
                  Axe : $g(z) = \frac{\sqrt2}{2}z$ équivaut à $i\bar{z} = z$, soit $y = x$ :
                  \fbox{$\Delta$ est la bissectrice de l'angle $\widehat{BAD}$} et $g = h\left(A,\frac{\sqrt2}{2}\right) \circ s_{\Delta}$.
      \end{enumerate}
\end{description}
\end{sol}
%%%%%%%%%%%%%%%%%%%%%%%%%%%%%%%%%%%%%%%%%%%%%%%%%%%%%%%%%%%%%%%%
\begin{sol}[Exercice~\ref{exo:sim-apollonius}]
\begin{enumerate}
      \item Rapport $\frac{OB}{OA} = \frac{\sqrt2}{2}$, angle $(\overrightarrow{OA},\overrightarrow{OB}) = \frac{\pi}{4}$ ;
            \fbox{$s : z' = \frac{1+i}{2}z$.}
      \item $\frac{1+i}{2}z_C = 2$, donc $z_C = \frac{4}{1+i} = 2(1 - i)$. \fbox{$C(2 - 2i)$.}
      \item $s(C) = A$ et $s(M) = M'$, et $s$ multiplie les distances par $\frac{\sqrt2}{2}$ : $AM' = \frac{\sqrt2}{2}CM$.
      \item \begin{enumerate}
                  \item $AM' = OM \iff \frac{\sqrt2}{2}CM = OM \iff CM^2 = 2OM^2$ (longueurs positives), soit $CM^2 - 2OM^2 = 0$.
                  \item $1 - 2 = -1 \neq 0$, donc $G$ existe : $z_G = \frac{z_C - 2z_O}{-1} = -2 + 2i$.
                        Pour tout $M$ : $CM^2 - 2OM^2 = (1 - 2)MG^2 + GC^2 - 2GO^2$ (en développant $\overrightarrow{MC} = \overrightarrow{MG} + \overrightarrow{GC}$, etc.,
                        le terme $2\overrightarrow{MG}\cdot(\overrightarrow{GC} - 2\overrightarrow{GO})$ est nul).
                        $GC^2 = |4 - 4i|^2 = 32$ et $GO^2 = |2 - 2i|^2 = 8$, d'où $CM^2 - 2OM^2 = -MG^2 + 16$.
                  \item $M \in \Gamma \iff MG^2 = 16 \iff MG = 4$ : \fbox{$\Gamma$ est le cercle de centre $G(-2+2i)$ et de rayon $4$.}
                        $|z_E - z_G| = |4| = 4$ : $E \in \Gamma$. (Directement : $s(E)$ a pour affixe $(1+i)^2 = 2i$, $AM' = |2i - 2| = 2\sqrt2 = OE$.)
            \end{enumerate}
\end{enumerate}
\end{sol}
%%%%%%%%%%%%%%%%%%%%%%%%%%%%%%%%%%%%%%%%%%%%%%%%%%%%%%%%%%%%%%%%
\begin{sol}[Exercice~\ref{exo:sim-triangle-equilateral}]
\begin{description}
      \item[I.]
      \begin{enumerate}
            \item Soit $I$ le milieu de $[BC]$ : $I$ est le barycentre de $(B;1)$, $(C;1)$, donc (associativité)
                  $G$ est le barycentre de $(A;2)$, $(I;2)$ : \fbox{$G$ est le milieu de $[AI]$}, où $[AI]$ est la médiane (et hauteur) issue de $A$.
            \item $\overrightarrow{MM'} = 2\overrightarrow{MA} + \overrightarrow{MB} + \overrightarrow{MC} = 4\overrightarrow{MG}$, donc
                  $\overrightarrow{GM'} = \overrightarrow{GM} + \overrightarrow{MM'} = \overrightarrow{GM} - 4\overrightarrow{GM} = -3\overrightarrow{GM}$.
                  \fbox{$f$ est l'homothétie de centre $G$ et de rapport $-3$.}
            \item On a $AI = \frac{a\sqrt3}{2}$, $GA = GI = \frac{a\sqrt3}{4}$ et $GB^2 = GC^2 = GI^2 + IB^2 = \frac{3a^2}{16} + \frac{a^2}{4} = \frac{7a^2}{16}$.
                  Par la réduction habituelle,
                  $2MA^2 + MB^2 + MC^2 = 4MG^2 + 2GA^2 + GB^2 + GC^2 = 4MG^2 + \frac{6a^2}{16} + \frac{14a^2}{16} = 4MG^2 + \frac{5a^2}{4}$.
                  La condition devient $4MG^2 = \frac{3a^2}{4}$, soit $MG = \frac{a\sqrt3}{4} = GA$.
                  \fbox{$\mathcal{E}$ est le cercle de centre $G$ passant par $A$} (il passe aussi par $I$).
            \item La somme des coefficients est nulle. En prenant $A$ comme origine :
                  $MB^2 + MC^2 - 2MA^2 = (\overrightarrow{MA} + \overrightarrow{AB})^2 + (\overrightarrow{MA} + \overrightarrow{AC})^2 - 2MA^2
                  = AB^2 + AC^2 + 2\overrightarrow{MA}\cdot(\overrightarrow{AB} + \overrightarrow{AC}) = 2a^2 + 4\,\overrightarrow{MA}\cdot\overrightarrow{AI}$.
                  La condition équivaut à $\overrightarrow{AM}\cdot\overrightarrow{AI} = 0$ :
                  \fbox{$\Delta$ est la droite passant par $A$ et perpendiculaire à $(AI)$, c'est-à-dire la parallèle à $(BC)$ menée par $A$.}
      \end{enumerate}
      \item[II.]
      \begin{enumerate}
            \item \begin{enumerate}
                        \item $BJ = IA = \frac{a\sqrt3}{2}$ (médianes du triangle équilatéral) et $\overrightarrow{BJ} \neq \overrightarrow{IA}$.
                              Il existe donc un unique déplacement envoyant $B$ sur $I$ et $J$ sur $A$, et ce n'est pas une translation :
                              c'est une rotation $r_1$.
                        \item $KB = KI = \frac{a}{2}$ et $KJ = KA = \frac{a}{2}$ (droites des milieux) : $K$ est sur les médiatrices de $[BI]$ et de $[JA]$,
                              qui sont distinctes ; c'est le centre de $r_1$.
                              Le triangle $KBI$ est équilatéral et direct (même orientation que $ABC$), donc $(\overrightarrow{KB},\overrightarrow{KI}) = \frac{\pi}{3}$.
                              \fbox{$r_1 = r\left(K,\frac{\pi}{3}\right)$.}
                  \end{enumerate}
            \item \begin{enumerate}
                        \item $r_2(J) = t(r_1(J)) = t(A) = J$. $r_2$ est la composée d'une rotation d'angle $\frac{\pi}{3}$ et d'une translation :
                              c'est une rotation d'angle $\frac{\pi}{3}$, de centre le point invariant $J$. \fbox{$r_2 = r\left(J,\frac{\pi}{3}\right)$.}
                        \item On cherche $\Delta_1$ passant par $K$ avec $(\Delta_1,(KC)) = \frac{\pi}{6}\ [\pi]$, et $\Delta_2$ passant par $J$ avec $(\Delta_2,(JC)) = \frac{\pi}{6}\ [\pi]$.
                              $\overrightarrow{KE} = 2\overrightarrow{KI}$ et $(\overrightarrow{AB},\overrightarrow{KI}) = \frac{\pi}{3}$, $(\overrightarrow{AB},\overrightarrow{KC}) = \frac{\pi}{2}$ :
                              donc $((KE),(KC)) = \frac{\pi}{6}$ et \fbox{$\Delta_1 = (KE)$.}
                              Le triangle $JIE$ est isocèle en $I$ ($IJ = IE = \frac{a}{2}$) avec $\widehat{JIE} = \frac{2\pi}{3}$, donc $\widehat{IJE} = \frac{\pi}{6}$ ;
                              comme $(JI) \parallel (AB)$ et $((AB),(JC)) = ((AB),(AC)) = \frac{\pi}{3}$, on obtient $((JE),(JC)) = \frac{\pi}{3} - \frac{\pi}{6} = \frac{\pi}{6}$ :
                              \fbox{$\Delta_2 = (JE)$.}
                        \item $t = r_2 \circ r_1^{-1}$ et $r_1^{-1} = s_{(KE)} \circ s_{(KC)}$, donc
                              $t \circ s_{(KC)} = r_2 \circ s_{(KE)} \circ s_{(KC)} \circ s_{(KC)} = r_2 \circ s_{(KE)} = s_{(JC)} \circ s_{(JE)} \circ s_{(KE)} = f$.
                        \item $s_{(KC)}$ échange $A$ et $B$ et envoie $I$ sur $J$ ; puis $t_{\overrightarrow{AJ}}$ : $f(B) = t(A) = J$, $f(I) = t(J) = C$,
                              $f(K) = t(K) = I$ (car $\overrightarrow{KI} = \overrightarrow{AJ}$). \fbox{L'image du triangle $BIK$ est le triangle $JCI$.}
                              $f$ est la composée de trois réflexions, donc un antidéplacement.
                              On décompose $\overrightarrow{AJ} = \frac14\overrightarrow{AB} + \frac12\overrightarrow{KC}$, avec $\frac14\overrightarrow{AB} \perp (KC)$ et
                              $\frac12\overrightarrow{KC} \parallel (KC)$. Alors $t_{\frac14\overrightarrow{AB}} \circ s_{(KC)} = s_{\Droite}$, où $\Droite$ est l'image de $(KC)$
                              par $t_{\frac18\overrightarrow{AB}}$ : la parallèle à $(KC)$ passant par le milieu de $[KI]$.
                              \fbox{$f = t_{\frac12\overrightarrow{KC}} \circ s_{\Droite}$ : symétrie glissée d'axe $\Droite$ et de vecteur $\frac12\overrightarrow{KC}$ non nul.}
            \end{enumerate}
      \end{enumerate}
      \item[III.] Ici $a = 4~\mathrm{cm}$ et l'unité vaut $2~\mathrm{cm}$ : $z_A = 0$, $z_B = 2$, $z_C = 1 + i\sqrt3$, $z_K = 1$,
            $z_I = \frac32 + \frac{\sqrt3}{2}i$, $z_J = \frac12 + \frac{\sqrt3}{2}i$, $z_E = 2z_I - z_K = 2 + i\sqrt3$, $z_D = \frac{z_A + z_B + z_C}{3} = 1 + \frac{\sqrt3}{3}i$.
      \begin{enumerate}
            \item \begin{enumerate}
                        \item $t_{\overrightarrow{AJ}} : z' = z + \frac12 + \frac{\sqrt3}{2}i$ ; $(KC)$ est la droite $x = 1$, donc $s_{(KC)} : z' = -\bar{z} + 2$.
                        \item $f : z' = -\bar{z} + \frac52 + \frac{\sqrt3}{2}i$, soit \fbox{$x' = -x + \frac52$, $y' = y + \frac{\sqrt3}{2}$.}
                              (On retrouve l'axe $x = \frac54$ et le vecteur $\left(0;\frac{\sqrt3}{2}\right) = \frac12\overrightarrow{KC}$.)
                  \end{enumerate}
            \item $a_s = \frac{z_D - z_I}{z_C - z_E} = \frac{-\frac12 - \frac{\sqrt3}{6}i}{-1} = \frac12 + \frac{\sqrt3}{6}i = \frac{\sqrt3}{3}e^{i\pi/6}$ et
                  $b_s = z_I - a_s z_E = 1 - \frac{\sqrt3}{3}i$.
                  \fbox{$s : z' = \left(\frac12 + \frac{\sqrt3}{6}i\right)z + 1 - \frac{\sqrt3}{3}i$} :
                  rapport $\frac{\sqrt3}{3}$, angle $\frac{\pi}{6}$, centre $\frac{b_s}{1 - a_s} = \frac{1 - \frac{\sqrt3}{3}i}{\frac12 - \frac{\sqrt3}{6}i} = 2$, c'est-à-dire \fbox{le point $B$.}
      \end{enumerate}
\end{description}
\end{sol}
%%%%%%%%%%%%%%%%%%%%%%%%%%%%%%%%%%%%%%%%%%%%%%%%%%%%%%%%%%%%%%%%
\begin{sol}[Exercice~\ref{exo:sim-triangle-isocele}]
On utilisera, pour contrôler, le repère de la partie B.II : $z_A = 0$, $z_B = 1$, $z_C = i$, $z_D = 2$, $z_O = 1 + \frac{i}{2}$ (unité $a$).
\begin{description}
      \item[Partie A.]
      \begin{enumerate}
            \item $-1 + 2 + 1 = 2 \neq 0$ ; $\overrightarrow{AG} = \frac12(2\overrightarrow{AB} + \overrightarrow{AC}) = \frac12(\overrightarrow{AD} + \overrightarrow{AC})$ :
                  \fbox{$G$ est le milieu de $[CD]$, c'est-à-dire $G = O$.}
            \item $-\overrightarrow{MA} + 2\overrightarrow{MB} + \overrightarrow{MC} = 2\overrightarrow{MO}$ ; la condition est $2MO = a$ :
                  \fbox{$\mathscr{C}$ est le cercle de centre $O$ et de rayon $\frac{a}{2}$.}
            \item $-MA^2 + 2MB^2 + MC^2 = 2MO^2 - OA^2 + 2OB^2 + OC^2$, avec $OA^2 = OC^2 = \frac{5a^2}{4}$ (rayon de $(\Gamma)$) et $OB^2 = \frac{a^2}{4}$ :
                  la somme vaut $2MO^2 + \frac{a^2}{2}$. La condition donne $MO^2 = \frac{a^2}{4}$ :
                  \fbox{$\mathscr{C}' = \mathscr{C}$, le cercle de centre $O$ et de rayon $\frac{a}{2}$.}
            \item La somme des coefficients est nulle, donc le vecteur ne dépend pas de $M$ ; avec $M = A$ :
                  $\overrightarrow{W} = 2\overrightarrow{AB} - \overrightarrow{AC} = \overrightarrow{AD} - \overrightarrow{AC}$. \fbox{$\overrightarrow{W} = \overrightarrow{CD}$.}
      \end{enumerate}
      \item[Partie B.]
      \textbf{I.}
      \begin{enumerate}
            \item $k = \frac{BC}{DB} = \frac{a\sqrt2}{a} = \sqrt2$ et
                  $\alpha = (\overrightarrow{DB},\overrightarrow{BC}) = (\overrightarrow{BA},\overrightarrow{BC}) = -\frac{\pi}{4}\ [2\pi]$
                  (le triangle $ABC$ est rectangle isocèle direct en $A$).
            \item $s \circ s$ est la similitude directe de centre $I$, de rapport $2$ et d'angle $-\frac{\pi}{2}$, et $s \circ s(D) = s(B) = C$.
                  Donc $(\overrightarrow{ID},\overrightarrow{IC}) = -\frac{\pi}{2}\ [2\pi]$ et $IC = 2ID$.
            \item D'après (\ref{eq:angle}), $\widehat{DIC}$ est droit : $I$ est sur le cercle de diamètre $[CD]$, c'est-à-dire $(\Gamma)$.
                  Par Pythagore, $ID^2 + IC^2 = CD^2 = AD^2 + AC^2 = 5a^2$ ; avec (\ref{eq:rapport}), $5ID^2 = 5a^2$, donc $ID = a$.
                  Enfin $s(I) = I$ et $s(D) = B$ donnent $IB = \sqrt2\, ID = a\sqrt2 = BC$.
            \item $OI = OC$ (rayons de $(\Gamma)$) et $BI = BC$ : $O$ et $B$ sont sur la médiatrice de $[IC]$, et $O \neq B$, donc $(OB)$ est cette médiatrice.
                  Le quadrilatère $CADI$ a ses côtés opposés de même longueur ($ID = CA = a$, $IC = AD = 2a$) ; $A$ et $I$ sont de part et d'autre de $(CD)$
                  (les angles $(\overrightarrow{AD},\overrightarrow{AC}) = \frac{\pi}{2}$ et $(\overrightarrow{ID},\overrightarrow{IC}) = -\frac{\pi}{2}$ sont opposés) :
                  c'est un parallélogramme ayant un angle droit, \fbox{$CADI$ est un rectangle}, et $I$ est le symétrique de $A$ par rapport à $O$.
      \end{enumerate}
      \textbf{II.}
      \begin{enumerate}
            \item $z_B = 1$, $z_C = i$, $z_D = 2$.
            \item $a_s = \frac{z_C - z_B}{z_B - z_D} = \frac{-1 + i}{-1} = 1 - i$, $b_s = z_B - a_s z_D = 1 - 2 + 2i = -1 + 2i$.
                  \fbox{$s : z' = (1 - i)z - 1 + 2i$} ; rapport $\sqrt2$, angle $-\frac{\pi}{4}$.
                  $z_0 = \frac{-1 + 2i}{1 - (1 - i)} = \frac{-1 + 2i}{i} = 2 + i$ : \fbox{$I(2 + i)$}, quatrième sommet du rectangle $CADI$ ($z_I = z_C + z_D - z_A$).
      \end{enumerate}
      \textbf{III.} $f$ fixe $A$, donc $f : z' = \lambda\bar{z}$ ; $f(C) = D$ donne $\lambda(-i) = 2$, $\lambda = 2i$.
      \fbox{$f : z' = 2i\,\bar{z}$} : similitude indirecte de centre $A$, de rapport $2$ ;
      son axe ($z' = 2z$, soit $i\bar{z} = z$) est la droite $y = x$, bissectrice de l'angle $\widehat{BAC}$.
      Forme réduite : $f = h(A,2) \circ s_{\Delta}$ avec $\Delta : y = x$.
\end{description}
\end{sol}
\end{corrige}
